\chapter{Definitions and conventions}\label{app:definitions}

The following conventions fix the object under discussion and the translations between its arithmetic, geometric and local descriptions. They are part of the claim. Changing a coefficient shift, the Vandermonde orientation or a Frobenius convention changes the formulas that a computation must verify.

\section{The coefficient and its local parameters}
The normalized form is
\[
\Delta(q)=q\prod_{m\ge1}(1-q^m)^{24}
         =\sum_{n\ge1}\tau(n)q^n,\qquad \tau(1)=1.
\]
All coefficient indices in this book are positive integers. A stored value at index zero is a programming sentinel unless a displayed formula explicitly introduces the coefficients of $\Delta/q$. For the latter coefficients, $f_m=\tau(m+1)$ and $f_0=1$.

At a prime $p$, the default local parameters are $a=\tau(p)$ and $b=p^{11}$. The scalar recurrence has $T_0=1$, $T_1=a$ and $T_r=aT_{r-1}-bT_{r-2}$ for $r\ge2$. Thus $T_r=\tau(p^r)$. A construction that permits an arbitrary integer $a$ states that change explicitly. Such a recurrence need not be the local sequence of $\Delta$.

The symbol $T_k(p)$ in the trace-formula chapter denotes a Hecke operator on the weight-$k$ cusp-form space. It is an operator, whereas $T_r$ in a scalar recurrence is an integer or an integer polynomial. The weight $k$ and the exponent $r$ serve different roles.

The valuation $v_\ell(x)$ is the exponent of $\ell$ in a nonzero rational number $x$, with $v_\ell(0)=\infty$. The quadratic character used in the radial extraction is
\[
\chifive(d)=
\begin{cases}
0&5\mid d,\\
1&d\equiv1,4\pmod5,\\
-1&d\equiv2,3\pmod5.
\end{cases}
\]
The character value and the tenth power in the scaling formula are both essential.

\section{Packets, content and the coefficient shift}
Coordinates have indices $0,1,2,3,4$. A shell is
\[
\Vn=\left\{v\in\Z^5:
v_i\equiv i\pmod5,\quad
\sum_i v_i=0,\quad
\sum_i v_i^2=10n\right\}.
\]
The standard orientation is
\[
D(v)=\prod_{0\le i<j\le4}(v_i-v_j),\qquad W(v)=D(v)/288.
\]
The identity is $\tau(n)=\sum_{v\in\Vn}W(v)$. No absolute values occur in that sum. Distinct residue classes ensure that no Vandermonde factor is zero. Integrality of $W$ is proved separately from the core identity.

The corresponding five-core partition has size $n-1$. In particular, the coefficient at $n=1$ corresponds to the empty partition. The involution is $J(v)_i=-v_{-i\bmod5}$. A $J$-orbit may contain one or two packets, and its mass is the sum of their weights.

The content of a packet is the positive gcd of its five coordinates. Content-one packets are primitive. Multiplication followed by the prescribed residue reindexing gives
\[
\tau(n)=\sum_{g^2\mid n}\chifive(g)g^{10}R(n/g^2),
\qquad
R(n)=\sum_{d^2\mid n}\mu(d)\chifive(d)d^{10}\tau(n/d^2).
\]
The boundary values are $R(1)=1$ and $R(p)=\tau(p)$. At a prime shell the only possible content is one, so the extraction has not created a new prime coefficient.

\section{Three polynomial operations}
The ordinary hyperplane Laplacian is $\Delta_H$. The conjugated operator $\mathcal L$ is characterized by $\Delta_H(DQ)=D\mathcal LQ$. Harmonicity of $DQ$ means $\mathcal LQ=0$. It does not mean that $Q$ itself is harmonic.

A universal null is a polynomial $Q$ for which
\[
\sum_{v\in\Vn}W(v)Q(v)=0
\]
for every positive $n$. A shell-null polynomial satisfies that equation only at a specified shell. A polynomial interpolating prescribed values on a shell need satisfy neither condition. The difference between these three operations is responsible for several of the obstruction results.

For a universal null, replacing $W(v)$ by $W(v)(1+Q(v))$ preserves the shell total. A sign conclusion requires real corrected weights and at least one strict nonzero contribution. There is no order relation on complex weights. The all-polynomial content theorem assumes a finite symmetric polynomial and makes no claim for arbitrary infinite series or nonsymmetric functions.

\section{Local and global conventions}
For $\ell\ne p$, the arithmetic-Frobenius convention in the integral $\ell$-adic representation gives the characteristic polynomial
\[
X^2-\tau(p)X+p^{11}.
\]
The cohomological conventions in the filtered-module discussion are stated there separately. An abstract trace-zero filtered module and a representation realizing it are local objects. A comparison identifying its normalized Frobenius trace with the actual coefficient $\tau(p)$ is additional mathematical data.

The fixed-modulus survivor theorem concerns the actual joint finite image, including the cyclotomic factor. It does not assume that its constituent images are independent. Its quantifiers are
\[
\text{for every fixed }M,\quad \text{there exist suitable actual primes }p.
\]
The prime may change with $M$. The theorem does not supply one prime for every modulus.

\chapter{The symmetric Laplacian in coordinates}\label{app:operator}

The universal null construction depends on an exact differential operator. This appendix derives that operator from the Euclidean metric, rather than treating its power-sum formulas as a black box. It also explains the Newton reductions used by the implementation and gives two low-degree harmonic checks. All polynomial identities are on the hyperplane $H=\{v\in\R^5:\sum_i v_i=0\}$, with its induced Euclidean metric. The modularity and valence contracts remain the ones stated in Chapters~\ref{ch:theta} and~\ref{ch:nulls}.

Write $s_m=\sum_{i=0}^4v_i^m$, with $s_0=5$ and $s_1=0$, and let $D=\prod_{i<j}(v_i-v_j)$. We distinguish the ordinary Laplacian $\Delta_H$ from the conjugated operator $\mathcal L$ defined by $\Delta_H(DQ)=D\mathcal LQ$. The factor $D$ changes radial coefficients substantially. Confusing the two operators would invalidate the harmonicity test.

\section{Two descriptions of the same metric}
The orthogonal projection from $\R^5$ to $H$ is
\begin{equation}\label{appop:projection}
P=I_5-\frac15\mathbf1\mathbf1^T,
\end{equation}
where $\mathbf1=(1,1,1,1,1)^T$. Indeed, subtracting the coordinate mean leaves a zero coordinate sum. The matrix is symmetric, satisfies $P^2=P$, and has image $H$. For an ambient polynomial $F$, its gradient along $H$ is $P\nabla F$. Its Laplacian along $H$ is
\begin{equation}\label{appop:ambientlap}
\Delta_H(F|_H)
=\left(\sum_{i=0}^4\partial_i^2F
-\frac15\left(\sum_{i=0}^4\partial_i\right)^2F\right)\Big|_H.
\end{equation}
Only tangential derivatives enter. Thus this expression is independent of the polynomial extension away from $H$.

For direct symbolic computation, use four coordinates $x_0,\ldots,x_3$ and substitute
\[
v=(x_0,x_1,x_2,x_3,-x_0-x_1-x_2-x_3).
\]
The coordinate tangent vectors are $b_a=e_a-e_4$. Their Gram matrix and its inverse are
\begin{equation}\label{appop:chartmetric}
G=I_4+\mathbf1\mathbf1^T,
\qquad G^{-1}=I_4-\frac15\mathbf1\mathbf1^T,
\end{equation}
now with a four-entry vector $\mathbf1$. Matrix multiplication verifies the inverse. The factor $1/5$ remains because the missing fifth coordinate is determined by the four chosen coordinates.

The chart Laplacian is consequently
\begin{equation}\label{appop:chartlap}
\Delta_H=\sum_{a,b=0}^3
\left(\delta_{ab}-\frac15\right)\frac{\partial^2}{\partial x_a\partial x_b}.
\end{equation}
The coordinates $x_a$ are not orthonormal. Using only the four diagonal second derivatives would use a different metric. As a normalization check, $s_2=x^TGx$, so its Hessian is $2G$ and
\[
\Delta_Hs_2=\operatorname{tr}(G^{-1}2G)=8.
\]
The rank of the Euclidean space is four, giving twice that rank. Later we obtain $\mathcal Ls_2=48$, an intentionally different number.

\section{Conjugating by the Vandermonde}
The Vandermonde is alternating under coordinate permutations. Since $\Delta_H$ commutes with those permutations, $\Delta_HD$ is alternating of degree eight. Every alternating polynomial vanishes on each hyperplane $v_i=v_j$. In the polynomial ring of $H$, the ten corresponding linear factors are distinct and divide it. Their product has degree ten. Therefore $\Delta_HD=0$.

On the open set where $D\ne0$, the product rule gives
\begin{equation}\label{appop:conjugation}
\frac{\Delta_H(DQ)}{D}
=\Delta_HQ+2\nabla_H\log D\cdot\nabla_HQ.
\end{equation}
Here $\nabla\log D$ denotes the rational logarithmic derivative $\nabla D/D$. It does not require choosing a complex logarithm or assuming that $D$ is positive. For symmetric $Q$, the apparent poles cancel and the right side extends to a polynomial on all of $H$.

There is a second explanation of the extension. The polynomial $\Delta_H(DQ)$ is alternating, so $D$ divides it. Its quotient is symmetric. Thus the operator $\mathcal L$ is well defined on symmetric polynomials even at coordinate collisions. The packet shell itself has no collisions because its residues are distinct. The polynomial extension is nevertheless needed when constructing and comparing symbolic identities.

Translation invariance of $D$ gives
\[
\sum_i\frac{\partial_iD}{D}=0,
\qquad
\frac{\partial_iD}{D}=\sum_{j\ne i}\frac1{v_i-v_j}.
\]
Its ambient logarithmic gradient is already tangential. This removes a potential extra projection term in the derivation below.

\section{Deriving the power-sum operator}
For $m\ge2$, differentiation and projection give
\begin{align}
(\nabla_Hs_m)_i
 &=m\left(v_i^{m-1}-\frac15s_{m-1}\right),\label{appop:gradsm}\\
\Delta_Hs_m
 &=\frac45m(m-1)s_{m-2}.\label{appop:lapsm}
\end{align}
In the logarithmic term, the coordinate-mean subtraction vanishes because the logarithmic gradient has zero sum. Pair the ordered indices $i,j$ and $j,i$ to obtain
\[
2\nabla_H\log D\cdot\nabla_Hs_m
=2m\sum_{i<j}\frac{v_i^{m-1}-v_j^{m-1}}{v_i-v_j}.
\]
The divided difference is the finite polynomial
\[
\frac{v_i^{m-1}-v_j^{m-1}}{v_i-v_j}
=\sum_{a=0}^{m-2}v_i^av_j^{m-2-a}.
\]
The sum over $a$ is symmetric in the two coordinates. Hence twice the sum over $i<j$ equals the sum over $i\ne j$. For each $a$, that ordered sum is $s_as_{m-2-a}-s_{m-2}$. There are $m-1$ values of $a$. We obtain
\[
2\nabla_H\log D\cdot\nabla_Hs_m
=m\sum_{a=0}^{m-2}s_as_{m-2-a}
-m(m-1)s_{m-2}.
\]
Adding \eqref{appop:lapsm} proves
\begin{equation}\label{appop:lsm}
\boxed{\mathcal Ls_m
=m\sum_{a=0}^{m-2}s_as_{m-2-a}
-\frac{m(m-1)}5s_{m-2}.}
\end{equation}
This formula is asserted for $m\ge2$. The boundary values are separately $\mathcal Ls_0=0$ and $\mathcal Ls_1=0$, since $s_0$ is constant and $s_1$ vanishes on $H$. No negative-index power sum is introduced.

The first values are
\begin{equation}\label{appop:firstoperator}
\mathcal Ls_2=48,\qquad
\mathcal Ls_3=0,\qquad
\mathcal Ls_4=\frac{188}{5}s_2,\qquad
\mathcal Ls_5=46s_3.
\end{equation}
For example, at $m=2$ the first term is $2s_0^2=50$, while the subtraction is $2s_0/5=2$. The additional $40$ compared with $\Delta_Hs_2=8$ comes from differentiating the degree-ten factor $D$.

\section{The chain-rule metric and polynomial assembly}
The same projected gradients give
\begin{equation}\label{appop:chainmetric}
\Gamma_{mn}:=\nabla_Hs_m\cdot\nabla_Hs_n
=mn\left(s_{m+n-2}-\frac15s_{m-1}s_{n-1}\right).
\end{equation}
Expanding the product in \eqref{appop:gradsm} gives the first power sum once and the coordinate-mean product with coefficient $-1/5$. This identity includes $m,n\ge1$, with the evident zero tangential gradient for $s_1$.

If $Q=F(s_2,s_3,s_4,s_5)$, the ordinary second-order chain rule and \eqref{appop:conjugation} yield
\begin{equation}\label{appop:chainrule}
\mathcal LQ
=\sum_{m=2}^5\frac{\partial F}{\partial s_m}\mathcal Ls_m
+\sum_{m,n=2}^5
\frac{\partial^2F}{\partial s_m\partial s_n}\Gamma_{mn}.
\end{equation}
The second sum runs over ordered pairs. Thus a mixed derivative contributes twice when $m\ne n$. If one instead sums over $m\le n$, the off-diagonal terms require a factor two. These are two equivalent conventions, but their factors must not be mixed.

Useful instances are
\begin{align}
\mathcal L(s_2^2)&=104s_2,\label{appop:lradialtwo}\\
\mathcal L(s_2^3)&=168s_2^2,\label{appop:lradialthree}\\
\mathcal L(s_2s_4)&=64s_4+\frac{188}{5}s_2^2,\label{appop:ls2s4}\\
\mathcal L(s_3^2)&=18s_4-\frac{18}{5}s_2^2.\label{appop:ls3square}
\end{align}
For the first identity, $2s_2\mathcal Ls_2+2\Gamma_{22}=96s_2+8s_2$. For the mixed product, $s_4\mathcal Ls_2+s_2\mathcal Ls_4+2\Gamma_{24}$ gives the displayed coefficients. Each calculation is a check of the ordered-pair convention.

The metric formula may produce $s_6,s_7,s_8$. They are not new independent variables. The invariant ring on $H$ is generated by $s_2,s_3,s_4,s_5$, and Newton's identities reduce higher indices exactly.

\section{Newton reductions with zero first power sum}
Let $e_j$ be the elementary symmetric polynomials in the five coordinates. Their characteristic polynomial is
\[
\prod_i(T-v_i)=T^5+e_2T^3-e_3T^2+e_4T-e_5,
\]
because $e_1=s_1=0$. The first Newton identities give
\begin{align}
e_2&=-\frac12s_2,& e_3&=\frac13s_3,\label{appop:e23}\\
e_4&=\frac18s_2^2-\frac14s_4,&
e_5&=\frac15s_5-\frac16s_2s_3.\label{appop:e45}
\end{align}
For instance, the degree-five identity is $s_5+e_2s_3-e_3s_2-5e_5=0$. The two middle terms combine to $-5s_2s_3/6$, fixing the sign in $e_5$.

Every coordinate satisfies its characteristic polynomial. Multiplying that equality by $v_i^{k-5}$ and summing, for $k\ge6$, gives
\begin{equation}\label{appop:newtonrec}
s_k=-e_2s_{k-2}+e_3s_{k-3}-e_4s_{k-4}+e_5s_{k-5}.
\end{equation}
This recurrence uses only lower indices. Substitution at $k=6,7$ yields
\begin{align}
s_6&=\frac34s_2s_4+\frac13s_3^2-\frac18s_2^3,\label{appop:s6}\\
s_7&=\frac7{10}s_2s_5+\frac7{12}s_3s_4
-\frac7{24}s_2^2s_3.\label{appop:s7}
\end{align}
The same rule gives
\begin{equation}\label{appop:s8}
s_8=\frac14s_4^2+\frac14s_2^2s_4
-\frac1{16}s_2^4+\frac8{15}s_3s_5.
\end{equation}
Each term has the required invariant degree. This supplies a useful local check before substitution into a larger operator expression.

At the first packet $v=(0,1,2,-2,-1)$, the first power sums are $s_2=10$, $s_3=0$, $s_4=34$, $s_5=0$. Formula \eqref{appop:s6} gives $255-125=130$, matching the direct sum of sixth powers. Formula \eqref{appop:s7} gives zero, matching cancellation of the opposite coordinates. Formula \eqref{appop:s8} gives $514$. These examples fix normalizations. The polynomial recurrences, rather than these examples, establish the identities on every vector in $H$.

\section{Radial factors and harmonic projection}
Let $h_\ell$ be homogeneous harmonic of degree $\ell$ on the rank-four space $H$. Euler's identity is $v\cdot\nabla_Hh_\ell=\ell h_\ell$. The product rule for the radial factor $s_2^j$ gives
\begin{equation}\label{appop:radial}
\Delta_H(s_2^jh_\ell)
=4j(\ell+j+1)s_2^{j-1}h_\ell,\qquad j\ge1.
\end{equation}
At $j=0$ the left side is zero. At $j=1,\ell=0$ the coefficient is $8$, agreeing with $\Delta_Hs_2$. At $j=1,\ell=10$ it is $48$, agreeing with $\Delta_H(s_2D)=48D$.

For a homogeneous polynomial $P$ of degree $M$, a concrete top-degree harmonic projection in dimension four is
\begin{equation}\label{appop:hprojection}
\mathcal H_M(P)=\sum_{j=0}^{\lfloor M/2\rfloor}
\frac{(-1)^j(M-j)!}{4^j j!M!}\,s_2^j\Delta_H^jP.
\end{equation}
For $M=0$ or $1$ the sum contains only $j=0$. In general, if $a_j$ denotes its coefficient, the product rule gives the cancellation condition
\[
a_j+4(j+1)(M-j)a_{j+1}=0.
\]
The factorial coefficients satisfy this condition. The last possible iterated Laplacian has no further nonzero derivative, so the resulting polynomial is harmonic. Its difference from $P$ is divisible by $s_2$.

Repeating the construction on that lower-degree quotient produces the harmonic decomposition
\begin{equation}\label{appop:decomposition}
P=\sum_{j=0}^{\lfloor M/2\rfloor}s_2^jh_{M-2j}.
\end{equation}
It is unique. If such a sum were zero, apply the Laplacian as many times as the largest occupied radial exponent. All smaller-exponent terms vanish. Formula \eqref{appop:radial} leaves a nonzero scalar times the last harmonic polynomial, so that polynomial must be zero. Descending removes every term.

When $P=DQ$ is alternating, each projected harmonic component is alternating because the Laplacian and multiplication by $s_2$ commute with permutations. Each nonzero component is divisible by $D$ and has degree at least ten. Thus for a correction of homogeneous degree $d$, only $j\le d/2$ can occur. This is the degree bound used in the quasimodular argument for nonhomogeneous content conservation.

Harmonic projection is not preservation of shell values. The elementary counterexample is $P=s_2D$, of degree $12$. Since $\Delta_HP=48D$ and the next Laplacian is zero,
\[
\mathcal H_{12}(s_2D)=s_2D-\frac1{48}s_2(48D)=0.
\]
Yet on $\Vn$ the original polynomial equals $10nD$, which is nonzero at every packet. The operation removes radial information. It cannot silently turn a shell interpolant into a universally null polynomial with the same values.

\section{Two explicit harmonic moment checks}
The degree-four correction
\[
K_4=130s_4-47s_2^2
\]
has
\[
\mathcal LK_4
=130\frac{188}{5}s_2-47(104s_2)=0.
\]
Thus $DK_4$ is harmonic. At the first packet, $W=1$ and $K_4=130\cdot34-47\cdot100=-280$. Its theta form has weight $16$, whose cusp space is spanned by $\Delta E_4$ with first coefficient one. Consequently the theta form is $-280\Delta E_4$, giving the moment identity already stated in Chapter~\ref{ch:nulls}.

Likewise, set
\[
K_6=90s_2s_4-27s_2^3-320s_3^2.
\]
Using \eqref{appop:lradialthree}--\eqref{appop:ls3square}, its operator image is
\begin{align*}
\mathcal LK_6={}&(5760-5760)s_4\\
&+(3384-4536+1152)s_2^2=0.
\end{align*}
At the first packet its value is $90\cdot10\cdot34-27\cdot1000=3600$. The weight-$18$ cusp space is spanned by $\Delta E_6$, again normalized with first coefficient one. Its theta form is therefore $3600\Delta E_6$.

These checks have different roles from the null certificates. Both polynomials are harmonic, but neither is universally null: their first coefficients are nonzero. Harmonicity supplies a modular-form space. The initial coefficients then identify the form or, when enough initial coefficients vanish, force it to be zero by valence. Skipping that second step would misidentify a nonzero moment as a permitted null correction.

The retained programs organize these same operations using sparse rational polynomials. A direct four-coordinate implementation can instead substitute $v_4=-\sum_{a=0}^3x_a$ into $DQ$ and apply \eqref{appop:chartlap}. Agreement between these organizations checks the symbolic construction through different algebraic representations. The universal justification remains the metric calculation, the exact polynomial identities and the stated harmonic theta and valence contracts. None of those statements alone excludes the unresolved zero prime parameter of $\Delta$.

\chapter{An exact computation workbook}
\label{app:computation-workbook}

This appendix gives calculations that a reader can reproduce without trusting a status flag. The objects are the normalized coefficients of $\Delta$, the auxiliary shell polynomials, and one rational transfer certificate at $31$. Each calculation ends at its actual mathematical boundary. Agreement of finite outputs does not supply the missing all-prime nonvanishing argument.

The complete integer inputs and outputs are in the companion evidence archive. Paths below are relative to its extracted project root. A JSON integer remains an integer even when it has many digits. An encoded rational pair $[a,b]$ means $a/b$, with a nonzero denominator. Converting these objects to floating point before checking an equality changes the contract.

\section{Three engines and one coefficient shift}
\label{sec:workbook-engines}

Write
\[
 F(q)=\frac{\Delta(q)}q=\prod_{k\ge1}(1-q^k)^{24}
      =\sum_{m\ge0}f_mq^m.
\]
Then $f_0=1$ and $\tau(n)=f_{n-1}$. This shift is essential. A program that stores $\tau(0)=0$ as a sentinel must not confuse that entry with the constant coefficient $f_0$.

For direct product expansion, terms of degree at most $m$ depend only on factors with $k\le m$. The exact binomial expansion
\[
 (1-q^k)^{24}=\sum_{j=0}^{24}(-1)^j\binom{24}{j}q^{kj}
\]
therefore supplies a finite algorithm for each desired coefficient. Three short calculations are
\begin{align*}
 f_1&=-24,\\
 f_2&=\binom{24}{2}-24=276-24=252,\\
 f_3&=-\binom{24}{3}+(-24)(-24)-24\\
    &=-2024+576-24=-1472.
\end{align*}
The last line comes respectively from the $q^3$ term in the first factor, the product of its $q$ term with the second factor's $q^2$ term, and the third factor's $q^3$ term. Thus it computes $\tau(4)$, not $\tau(3)$.

The logarithmic derivative provides a different exact engine. With $\sigma_1(j)=\sum_{d\mid j}d$, comparison of coefficients in
\[
 q\frac{F'(q)}{F(q)}=-24\sum_{j\ge1}\sigma_1(j)q^j
\]
gives
\begin{equation}
 m f_m=-24\sum_{j=1}^{m}\sigma_1(j)f_{m-j}
 \qquad(m\ge1).
 \label{eq:workbook-divisor-recurrence}
\end{equation}
For $m=3$, the divisor sums are $1,3,4$. Hence
\[
 f_3=-8\bigl(252+3(-24)+4\bigr)=-8\cdot184=-1472.
\]
An implementation should check that the numerator in \eqref{eq:workbook-divisor-recurrence} is divisible by $m$, rather than rounding the quotient.

Jacobi gives a third engine. Put $T_j=j(j+1)/2$ and $a_j=(-1)^j(2j+1)$. The series
\[
 A(q)=\prod_{k\ge1}(1-q^k)^3=1+\sum_{j\ge1}a_jq^{T_j}
\]
satisfies $F=A^8$. The identity $AqF'=8FqA'$ yields
\begin{equation}
 m f_m=\sum_{\substack{j\ge1\\T_j\le m}}
          a_j(9T_j-m)f_{m-T_j}.
 \label{eq:workbook-triangular-recurrence}
\end{equation}
At $m=3$, only $T_1=1$, $T_2=3$ contribute. Its numerator is
\[
 (-3)(9-3)(252)+5(27-3)(1)=-4536+120=-4416.
\]
Division by $3$ again gives $-1472$. The direct product, divisor recurrence and triangular recurrence reach the same integer through different algebraic operations.

\Needspace{11\baselineskip}
\begin{center}
\begin{tabular}{r r r r}
\toprule
$m$ & $f_m=\tau(m+1)$ & $m$ & $f_m=\tau(m+1)$\\
\midrule
0 & $1$ & 5 & $-6048$\\
1 & $-24$ & 6 & $-16744$\\
2 & $252$ & 7 & $84480$\\
3 & $-1472$ & 8 & $-113643$\\
4 & $4830$ & 9 & $-115920$\\
\bottomrule
\end{tabular}
\end{center}

The current independent comparison covers every $\tau(n)$ for $1\le n\le512$. The historical final-arithmetic run's direct-product comparison stopped at $128$, while the packet-tree run compared through $512$. These are different documented ranges. The new three-engine table supplies its own bound and does not infer one from an older summary.

\section{Reconstructing the auxiliary polynomial}
\label{sec:workbook-polynomial}

For a packet $v$, set $c_i=(v_i-i)/5$ and form
\[
 P_v(T)=\prod_{0\le i<j\le4}\bigl((i-j)+T(c_i-c_j)\bigr).
\]
Its physical value is $P_v(5)=288W(v)$. Summing the full shell gives
\[
 P_n(T)=\sum_{v\in\Vn}P_v(T)=\sum_{j=0}^{10}B_jT^j,
 \qquad P_n(5)=288\tau(n).
\]
The polynomial has a different variable from the prime index. A root in $T$ or in a rescaled variable does not represent a nearby prime coefficient.

For a nonzero $P_n$, use the convention $v_5(0)=+\infty$ and define
\[
 \mu_n=\min_j\bigl(j+v_5(B_j)\bigr),\qquad
 G_n(U)=5^{-\mu_n}P_n(5U)=\sum_{j=0}^{10}C_jU^j.
\]
Every $C_j=B_j5^{j-\mu_n}$ is an integer. The power of $5$ should be removed only after verifying divisibility for each coefficient. Since $288$ is a $5$-adic unit,
\begin{equation}
 v_5(\tau(n))=\mu_n+v_5(G_n(1)).
 \label{eq:workbook-normalized-valuation}
\end{equation}
All coefficients in the next two tables are exact integers from the freshly reproduced shell data. They allow a reader to check the normalization and every displayed residue without re-enumerating a large shell.

\section{The simple branch at 499}
\label{sec:workbook-499}

At $n=499$ the shell has $500$ packets. Of their weights, $262$ are positive and $238$ negative. Their sum is
\[
 \tau(499)=-108877719272500.
\]
The minimum layer valuation is $\mu_{499}=3$.

\Needspace{17\baselineskip}
\begin{center}
\begingroup\small
\setlength{\tabcolsep}{5pt}
\begin{tabular}{r r r r}
\toprule
$j$ & $B_j$ in $P_{499}$ & $C_j$ in $G_{499}$ & $j+v_5(B_j)$\\
\midrule
0 & $144000$ & $1152$ & 3\\
1 & $-360000$ & $-14400$ & 5\\
2 & $1063040$ & $212608$ & 3\\
3 & $28418240$ & $28418240$ & 4\\
4 & $-148872808$ & $-744364040$ & 4\\
5 & $209474504$ & $5236862600$ & 5\\
6 & $2422451240$ & $302806405000$ & 7\\
7 & $16445000$ & $10278125000$ & 11\\
8 & $8107402848$ & $25335633900000$ & 8\\
9 & $-11893316304$ & $-185833067250000$ & 9\\
10 & $-1160632800$ & $-90674437500000$ & 12\\
\bottomrule
\end{tabular}
\endgroup
\end{center}

Summing the second coefficient column and using the normalization checks gives
\begin{align*}
 P_{499}(5)&=-31356783150480000=288\tau(499),\\
 G_{499}(1)&=-250854265203840.
\end{align*}
Reducing the normalized coefficients gives
\begin{align*}
 G_{499}(U)&\equiv2+3U^2\pmod5,\\
 G_{499}(U)&\equiv2+8U^2+15U^3+10U^4\pmod{25}.
\end{align*}
The first polynomial has the simple roots $1$ and $4$ modulo $5$. At $U=1$, the exact sum is $10$ modulo $25$ and its derivative is $1$ modulo $5$. Consequently
\[
 G_{499}(1+5h)\equiv10+5h\pmod{25}.
\]
The five possibilities are $10,15,20,0,5$ as $h$ runs from $0$ to $4$. The root in the branch $U\equiv1\pmod5$ therefore lifts to $U=16\pmod{25}$.

This lifted root is not the physical target. The target remains $U=1$, where $v_5(G_{499}(1))=1$. Formula \eqref{eq:workbook-normalized-valuation} gives $v_5(\tau(499))=3+1=4$.

The following exact root history records both simple branches. Each entry lies between $0$ and $5^k-1$ and evaluates to zero modulo $5^k$. Consecutive entries agree modulo the preceding power. These conditions are stronger than displaying an approximate root.

\Needspace{18\baselineskip}
\begin{center}
\begin{tabular}{r r r r}
\toprule
$k$ & $5^k$ & Root reducing to $1$ & Root reducing to $4$\\
\midrule
1 & 5 & 1 & 4\\
2 & 25 & 16 & 4\\
3 & 125 & 41 & 54\\
4 & 625 & 541 & 554\\
5 & 3125 & 3041 & 2429\\
6 & 15625 & 9291 & 14929\\
7 & 78125 & 9291 & 77429\\
8 & 390625 & 165541 & 233679\\
9 & 1953125 & 1337416 & 1014929\\
10 & 9765625 & 7196791 & 1014929\\
11 & 48828125 & 7196791 & 20546179\\
12 & 244140625 & 153681166 & 215858679\\
\bottomrule
\end{tabular}
\end{center}

The fixed degree $10$ does not impose a universal ceiling on the valuation of polynomial values at varying inputs. A simple $5$-adic root supplies arbitrarily accurate integer approximations. That observation concerns the auxiliary variable. It neither changes the value at $1$ nor supplies interpolation of the prime-indexed family.

\section{The double branch at 16249}
\label{sec:workbook-16249}

At $16249$ the shell contains $16250$ packets, with $8284$ positive and $7966$ negative weights. The exact coefficient is
\[
 \tau(16249)=133116215274958169206250.
\]
Here $\mu_{16249}=4$. The full coefficient table is as follows.

\Needspace{17\baselineskip}
\begin{center}
\begingroup\small
\setlength{\tabcolsep}{4pt}
\begin{tabular}{r r r r}
\toprule
$j$ & $B_j$ in $P_{16249}$ & $C_j$ in $G_{16249}$ & $j+v_5(B_j)$\\
\midrule
0 & $4680000$ & $7488$ & 4\\
1 & $-13320000$ & $-106560$ & 5\\
2 & $-199541600$ & $-7981664$ & 4\\
3 & $-26375464320$ & $-5275092864$ & 4\\
4 & $737398074300$ & $737398074300$ & 6\\
5 & $-39838758304992$ & $-199193791524960$ & 5\\
6 & $-746041848478320$ & $-18651046211958000$ & 7\\
7 & $-6990990388478280$ & $-873873798559785000$ & 8\\
8 & $21389931029028020$ & $13368706893142512500$ & 9\\
9 & $-64003552771660728$ & $-200011102411439775000$ & 9\\
10 & $3937759175521630080$ & $61527487117525470000000$ & 11\\
\bottomrule
\end{tabular}
\endgroup
\end{center}

The exact sums are
\begin{align*}
 P_{16249}(5)&=38337469999187952731400000,\\
 G_{16249}(1)&=61339951998700724370240.
\end{align*}
The first equals $288\tau(16249)$, and the second equals the first divided by $625$. Modulo $5$, the normalized polynomial is
\[
 3+U^2+U^3=(U-1)^2(U-2).
\]
Modulo $25$, the table gives
\[
 G_{16249}(U)\equiv13+15U+11U^2+11U^3+15U^5.
\]
Thus $G_{16249}(1)\equiv15\pmod{25}$ while $G'_{16249}(1)\equiv0\pmod5$. For every integer or $5$-adic integer $h$, polynomial Taylor expansion gives
\[
 G_{16249}(1+5h)\equiv G_{16249}(1)+5hG'_{16249}(1)
                    \equiv15\pmod{25}.
\]
Terms of order at least two have a factor $25$. This is a uniform obstruction throughout that residue branch, not merely the observation that five trial inputs fail. Every input congruent to $1$ modulo $5$ has normalized valuation exactly $1$.

The other residual root should not be discarded. At $U=2$, the normalized value is $5$ modulo $25$ and its derivative is $1$ modulo $5$. That simple branch lifts to $22$ modulo $25$. The conclusion is that the double branch at $1$ does not lift, rather than that the polynomial has no $5$-adic roots anywhere.

At the physical input $1$, the coefficient valuation is $4+1=5$. The two examples have different root geometry but both settle their actual target value by an exact modulo-$25$ calculation.

\section{A first failure needs a domain}
\label{sec:workbook-valuation-domain}

The tested auxiliary equality was
\[
 v_\ell(\tau(p))=v_\ell(p+1),
 \qquad p\ne\ell,\quad\ell\mid(p+1),\quad p\text{ prime}.
\]
The complete search through $150000$ contains $13848$ primes. The following counts refer to the indicated candidate sets, not to every prime without a residue condition.

\Needspace{11\baselineskip}
\begin{center}
\begingroup\small
\begin{tabular}{r r r r r r}
\toprule
$\ell$ & First $p$ & $v_\ell(\tau(p))$ & $v_\ell(p+1)$ & Candidates & Discrepancies\\
\midrule
2 & 81919 & 20 & 14 & 13847 & 1\\
3 & 8747 & 8 & 7 & 6941 & 5\\
5 & 499 & 4 & 3 & 3451 & 50\\
\bottomrule
\end{tabular}
\endgroup
\end{center}

The exact decisive values, including the discrepancy in the opposite $5$-adic direction, are
\begin{align*}
 \tau(499)&=-108877719272500,\\
 \tau(1249)&=27757672778617250,\\
 \tau(8747)&=348766170858826454196,\\
 \tau(81919)&=-904107824789465028682055680.
\end{align*}
For example, $1250=2\cdot5^4$, while the second integer has $5$-adic valuation $3$. Hence a proposed one-sided lift $v_5(\tau(p))\ge v_5(p+1)$ also fails at this actual prime.

These four coefficients were checked independently by the Eisenstein convolution. If $\sigma_r(n)=\sum_{d\mid n}d^r$, the coefficient of $(E_4E_8-E_6^2)/1728$ is obtained from
\begin{align*}
1728\tau(n)={}&240\sigma_3(n)+480\sigma_7(n)+1008\sigma_5(n)\\
 &+115200\sum_{k=1}^{n-1}\sigma_3(k)\sigma_7(n-k)\\
 &-254016\sum_{k=1}^{n-1}\sigma_5(k)\sigma_5(n-k).
\end{align*}
This computes the selected integers without the triangular recurrence. The numerator is divided exactly by $1728$. The replay checks this independent formula for $1\le n\le128$ and at the four displayed primes, giving $132$ checks.

The phrase ``first failure'' has two limits. It concerns the ordered search in its declared candidate set, and its exhaustive range ends at $150000$. A larger prime may supply another failure. The displayed nonzero integers refute an auxiliary valuation identity, not Lehmer's conjecture.

\section{Reading the rational transfer at 31}
\label{sec:workbook-transfer}

The degree-at-most-$22$ certificate has six coarse sectors and sixteen fine sectors. Both correction matrices have eighteen columns. Let
\[
 A\in\Z^{6\times16},\quad
 M_c\in\Z^{6\times18},\quad M_f\in\Z^{16\times18}.
\]
The fine-to-coarse parent map completely specifies $A$: its column $j$ has a $1$ in the parent row and zero elsewhere. The original fine masses are given below. Let $s_f$ denote their sixteen-entry column vector. Indices start at $0$.

\Needspace{15\baselineskip}
\begin{center}
\begin{tabular}{r r r @{\qquad} r r r}
\toprule
$j$ & Parent & Fine mass & $j$ & Parent & Fine mass\\
\midrule
0 & 0 & $2354352$ & 8 & 0 & $-11844448$\\
1 & 2 & $-19467448$ & 9 & 1 & $-7162848$\\
2 & 4 & $-832048$ & 10 & 0 & $-31156048$\\
3 & 4 & $15074752$ & 11 & 2 & $4356352$\\
4 & 4 & $1881152$ & 12 & 2 & $32821152$\\
5 & 5 & $4248552$ & 13 & 2 & $1262352$\\
6 & 5 & $7759752$ & 14 & 0 & $-7455448$\\
7 & 4 & $-31808448$ & 15 & 3 & $-12874848$\\
\bottomrule
\end{tabular}
\end{center}

For instance, row $0$ aggregates fine indices $0,8,10,14$:
\[
 2354352-11844448-31156048-7455448=-48101592.
\]
The six coarse masses are therefore
\[
 s_c=(-48101592,-7162848,18972408,
      -12874848,-15684592,12008304)^t.
\]
Their sum, and the sum of the fine masses, is $-52843168$. These original masses differ in order from the lexicographic orbit bars because the certificate orders the residue-sector keys.

The full machine certificate is
\begin{center}\small
\path{evidence/computation/replays/transport/p31_transport_certificate.json}.
\end{center}
It supplies an $18\times13$ rational matrix $N$ whose columns form a basis of $\ker M_c$, and the detail matrix $B=M_fN$ of size $16\times13$. The basic exact checks are
\[
 AM_f=M_c,\quad As_f=s_c,\quad M_cN=0,
 \quad B=M_fN,\quad AB=0.
\]
The last equation shows that a detail correction changes no coarse mass.

The rational row-reduction results are
\[
 \rank M_c=5,\qquad\rank M_f=15,\qquad\rank B=10.
\]
An independent check obtains the same lower bounds over both prime fields $\mathbb F_{65537}$ and $\mathbb F_{1000003}$. Every denominator used in a reduction is invertible there. A nonzero minor in either field proves a nonzero rational minor after clearing denominators.

The corresponding upper bounds have direct structural explanations. Each correction matrix has column sum zero, giving bounds $5$ and $15$. Every row of $A$ has a nonempty, disjoint support, so $\rank A=6$ and $\dim\ker A=10$. Since $AB=0$, the detail rank is at most $10$. The upper and lower bounds establish the displayed ranks exactly.

The coarse coefficient vector $a\in\Q^{18}$ makes
\[
 x=s_f+M_fa,\qquad Ax=(0,0,0,0,0,-52843168)^t.
\]
The certificate then chooses a coefficient vector $d\in\ker M_c$. Its final fine vector $y=x+M_fd$ is zero except at index $6$, where it equals $-52843168$.

One parent gives a concrete view of this second stage. Parent $5$ has fine children $5,6$. Just after the coarse stage their values are
\begin{align*}
 x_5&=-\frac{161598180629860124657534386744}
                 {23300823161076013875625},\\
 x_6&=-\frac{1069691132209170737342448593256}
                 {23300823161076013875625}.
\end{align*}
Their exact sum is $-52843168$. The detail stage adds $-x_5$ to child $5$ and $x_5$ to child $6$, giving $(0,-52843168)$ in that parent. All other parents retain their zero totals. This is a rational conservation calculation, not a numerical near-zero test.

To verify the whole certificate, parse the rational pairs, check the matrix dimensions, verify all five identities, and compare both stored stages entry by entry. Then check $M_cd=0$ and the single surviving final entry. The script \path{code/check_transport_certificate.py} performs the independent equality and prime-field checks.

This example does not assume a uniform rank theorem. Nor does full rank prove a uniformly nonzero survivor. At a shell whose original total were zero, complete zero-total transfer could still leave every final mass zero. At $31$ the separately computed exact total is nonzero, so the displayed survivor is authenticated for this one shell.

\section{Reproducing a run and interpreting its record}
\label{sec:workbook-manifests}

From the extracted source-set root, a new independent coefficient run is requested with
\begin{lstlisting}[basicstyle=\ttfamily\small,breaklines=true,frame=single,rulecolor=\color{Soft}]
python3 code/run_replays.py tau_independent \
  --output-root evidence/computation/reader_coefficients
\end{lstlisting}
The destination must be fresh. The launcher refuses to replace an existing run directory. Omitting the suite name requests the complete declared sequence of suites. For that complete reproduction, use a fresh destination such as \path{evidence/computation/reproduced}. The large modular comparison consumes the newly generated dyadic result in that same destination, so these dependent suites belong in their declared order.

A version $2$ manifest records the actual argument vector, exact arithmetic domain, semantic bounds, software versions and executed limits. Its \texttt{input\_artifacts} contain input checksums before and after execution. Its \texttt{declared\_results} name the scientific files the child must create. Output checksums cover those files and the execution logs.

Three questions remain separate. Did execution finish within its limits with unchanged inputs and all required outputs? Does the code implement the intended finite assertion? Does a proof show that the finite assertion decides the universal claim? An exit code of zero answers only part of the first question.

The count $704$ means $22$ selected modulus examples times $32$ exponents, not $704$ distinct primes. Of $24$ selected moduli, $11$ and $19$ had no observed matching prime through $10000$. That finite absence does not contradict positive density. The $24576$ synthetic cases mean $256$ moduli times three trace parameters times $32$ exponents. Their parameter $p=5M-1$ need not be prime, and their trace is not asserted to equal $\tau(p)$.

A checksum pins bytes. Run provenance records execution. A timeout records an execution limit. A reproduced obstruction can refute its stated auxiliary implication while leaving Lehmer unchanged. Attach each exact calculation to a claim of the correct scope.


\chapter{External theorem contracts}\label{app:sources}

This edition checks the specific theorem statements used in the argument. The check is targeted, not an exhaustive review of the literature or a certification of novelty. The sixteen source contracts below are recorded with versions, locators, access outcomes and limitations in the companion evidence. The bibliography supplies the corresponding primary or author-hosted references.

The source-dependent input and the project argument should be read together. For example, a vector-valued harmonic theta theorem does not by itself prove the scalar sign-line descent used here. That descent is developed in the main text. Similarly, a trace formula computes an exact difference, but its existence does not exclude a zero difference.

\section{Classical coefficient structure}
\paragraph{S01: Lehmer's original reduction.}
Lehmer's 1947 paper is the historical reference for the least-zero prime reduction and the multiplicative structure \cite{Lehmer}. The original scan was not text-readable through the checked access route, so this edition does not claim to have freshly reconstructed its original proof. The necessary coefficient identities are separately checked in readable modern primary sources, and the prime-power zero criterion is proved in the main text.

\paragraph{S06: Deligne's bound and normalization.}
The modular representation and purity inputs give $|\tau(p)|\le2p^{11/2}$ \cite{DeligneModular,DeligneWeil}. The modular paper's representation construction and the Weil theorem supply different parts of this implication. Their Frobenius conventions require translation. This upper bound includes zero and is never used as a positive lower bound.

\paragraph{S07: classical congruences.}
Serre's 1967/68 exposition supplies the checked modulo-$691$ identity and the modulo-$23$ prime classification \cite{SerreCongruences}. The latter excludes $p=23$. Representation existence is still conjectural in the relevant passage of that early exposition. The existence input in this book therefore uses the later established representation statement.

\paragraph{S10: a cited historical computation.}
Derickx, van Hoeij and Zeng establish the strict historical range
\[
n<816212624008487344127999
\]
in their checked 2014 version, Corollary~1.2 \cite{DHZ}. The endpoint is not asserted to be a zero. This source-certified range was not rerun for the present edition and is not presented as the latest record.

\paragraph{S11: scoped contemporary context.}
The checked 2026 version supplies a contemporary formulation and readable classical recurrence and parity statements \cite{TauContext2026}. Only those classical introductory inputs are used. The 2025 submission titled as a proof was withdrawn \cite{Withdrawn2025}. Its title is not accepted as evidence of a theorem. No exhaustive current-status or priority claim follows from this targeted check.

\paragraph{S16: readable classical leaves.}
Balakrishnan, Craig and Ono provide the modulo-$5$ congruence at equation~(1.2), the multiplicativity and recurrence statements at Theorem~2.2, and the direct odd-square parity argument in Section~3 of the checked 2020 version \cite{BCO}. These are the exact leaves used here. This edition has not audited the paper's separate integral-point or omitted-value arguments.

\section{Convolutions and trace moments}
\paragraph{S12: the Hurwitz moment normalization.}
Lygeros and Rozier's 2010 primary paper gives the trace moments and Catalan inversion used in Chapter~5 \cite{LRTrace}. The class-number function includes imprimitive forms with its specified weights. The $t=0$ term has a half-weight in the zeroth moment. The integral normalization is $6H$. Removing any of these qualifications would change the rows of the inverse system.

\paragraph{S14: Niebur's convolution.}
The universal formula is checked in the primary paper of Moriya and Smyth, Section~7, rather than inferred from finite coefficient agreement \cite{MoriyaSmyth}. The original 1975 Niebur proof was inaccessible through the checked route \cite{Niebur}. The attribution is retained with that access limitation. The size and divisibility alternatives remain compatible with a prime zero.

\paragraph{S15: differential equations.}
Fonseca records the normalized Ramanujan differential relations used in the polynomial computation \cite{Fonseca}. The project-specific $S_4$ cancellation is verified algebraically in the main text and by an exact symbolic remainder. Modularity and the level-one dimension input for the cusp identity are separate contracts \cite{Stein}. A finite output check is not substituted for the universal algebraic identity.

\section{Cores, harmonic theta series and valence}
\paragraph{S02: Han's packet identity.}
Han's Definition~1.1 and Theorems~1.1--1.2 in the checked first arXiv version provide the five-core specialization and Vandermonde product \cite{Han}. The specialization $t=5$, $z=25$ gives the shell norm $10n$ after the coefficient shift. With the orientation used here, the normalization is $D/288$. The identity permits cancellation and does not imply that its terms have one sign.

\paragraph{S13: packet cardinality.}
Garvan, Kim and Stanton's Theorem~4 and equation~(5.3) give the five-core count \cite{GKS}. In the present conventions this reads
\[
a_5(n-1)=\sum_{d\mid n}\chifive(d)\frac nd.
\]
For a prime different from $5$, the shell cardinality is $p+\chifive(p)$. This is the number of packets, not their weighted sum.

\paragraph{S03: harmonic theta modularity.}
Ibukiyama and Williams, Section~3.4 of the checked version, give the vector-valued harmonic theta theorem with weight $\nu+r/2$ \cite{IW}. Here the rank is four and the harmonic polynomial is alternating of degree $10+d$. The scalar full-level transformation follows from the separate sign-line descent in Chapter~8. It is not silently imported as part of the general theorem.

\paragraph{S04: the valence step.}
Stein's author exposition supplies the level-one valence bound and the relevant cusp-form dimensions \cite{Stein}. The vanishing test is applied only after modularity, holomorphy and the weight are established. Vanishing of sufficiently many initial Fourier coefficients then proves that specific modular form is zero. It does not establish an arbitrary polynomial's universal nullness without the modular prerequisites.

\section{Global congruences and local realization}
\paragraph{S05: the actual compatible representations.}
Serre's 1973 exposition gives the established integral rank-two representations and the arithmetic-Frobenius polynomial \cite{SerreBourbaki}. The fixed-modulus proof uses their actual joint finite quotient. It does not require an image classification, surjectivity onto an ambient matrix group, or independence across primes $\ell$.

\paragraph{S08: Chebotarev's theorem.}
The checked Stevenhagen--Lenstra exposition states the nonempty conjugacy-class theorem and distinguishes natural density \cite{Chebotarev}. In a finite Galois quotient the density is the class size divided by the group size. Complex conjugation supplies the class in the actual image. The theorem is applied separately for each fixed modulus. Nothing in the density formula interchanges the quantified modulus and prime.

\paragraph{S09: weak admissibility and admissibility.}
The Colmez--Fontaine author manuscript supplies the implication from weak admissibility to admissibility and the $N=0$ crystalline realization in the specified setting \cite{CF}. The constructed rank-two module has equal Hodge and Newton totals and no proper Frobenius-stable rational line. Those checks concern a local object. They do not identify it with the global representation of $\Delta$.

\section{Access and attribution limits}
The companion source records retain the exact checked versions and successful alternative author-hosted access routes. The two original-proof access limits above are explicit. The book does not label a cached excerpt, a source title, a search result or a historical status field as a completed proof audit.

The internal packet transport, nonhomogeneous content argument, interpolation analysis and arithmetic model constructions are reviewed as written arguments. Their correctness is distinct from whether any formulation has appeared elsewhere. This edition makes no historical-priority claim for them and no exhaustive novelty claim.

\chapter{Reproduction and evidence boundaries}\label{app:reproduction}

The companion archive is a reproducible record of this edition's finite claims. It preserves six recovered research ZIPs and all their original extracted file members. New executions write to separate directories. The present prose and audits determine what this reader claims. Historical labels remain available for comparison without being treated as current validation.

\section{What was rerun}
Sixteen fresh suites completed. Their version-two execution manifests validate current inputs, scientific results and logs. Thirteen historical scientific outputs agree exactly with the fresh results after removing only elapsed-time metadata fields named \texttt{runtime}, \texttt{runtime\_seconds} and \texttt{seconds}. Three additional suites provide independent coefficient, transfer and large-residue checks.

The most direct coefficient check covers every $1\le n\le512$. The triangular recurrence, the defining binomial product and an independent logarithmic-derivative divisor-sum engine agree exactly. All three use integers. This range is distinct from the historical final-arithmetic program's product check through $128$, and from the independent product check in the packet-tree program through $512$.

The arithmetic survivor search uses $\tau$ through $10000$ and twenty-four selected moduli. Twenty-two moduli have a matching prime in that bounded search, supplying $22\times32=704$ local recurrence checks. Moduli $11$ and $19$ have no observed witness in that search. This is a search outcome, not a contradiction of the density theorem. Modulus one is included as a trivial case.

The synthetic recurrence suite has $256$ moduli, three parameter choices and thirty-two positive exponents, for $24576$ checks. Its value $p=5M-1$ need not be prime. Separately, the $1229$ prime-indexed companion matrices through $10000$ are constructed local models. They are not sampled global Frobenius matrices.

\section{Packet and polynomial checks}
The five-core suite checks $405$ shell sums and $326585$ packet polynomials, with $308$ prime shell counts. Twenty-five coefficients also agree with an independent hook-partition calculation. The tree suite enumerates all shells through $512$ and the two additional shells $737$ and $2411$, giving $514$ shell sums and $95937$ packets. It checks $3084$ sector partitions and explicit sharp-boundary witnesses at sixty-four depths.

The shell at $737$ is a sparse additional witness. The replay does not enumerate every shell from $513$ through $736$. The universal threshold formula rests on the separate lower-bound proof and the symbolic witness families.

The retained universal-null suite checks fourteen exact rational sign certificates. Its degree-at-most-sixteen dual at $31$ concerns the seven stored corrections, not all possible corrections. The primitive transport suite checks thirty-six symbolic harmonic certificates, $128$ Mobius identities, $144$ profiles and the stored rational transfer. The independent $31$ calculation obtains coarse, fine and detail ranks $5$, $15$ and $10$ over the rationals and the two prime fields $65537$ and $1000003$, with structural upper bounds.

The content obstruction is the exact shell $121$. Its masses are
\[
-25439104668+25937424601=498319933.
\]
The computation establishes these integers. The theorem that every finite symmetric universal null conserves each content is a separate proof. Likewise, the interpolation suite's finite orbit checks and two explicit failures do not prove an all-degree fullness theorem.

\section{Large modular coefficients}
The valuation suite regenerates coefficients at all $13848$ primes through $150000$. Its proposed valuation equalities are tested only on their stated candidate residue classes. Independent Eisenstein convolutions verify the selected counterexample coefficients. The dyadic obstruction uses twelve selected primes through $1032191$, with nineteen fresh modular comparisons.

The CRT/NTT computation returns residues modulo $2^{19}$. It is not a reconstruction of those large integer tau coefficients. Within each convolution, however, the integer coefficient is reconstructed before reduction. With
\[
L=1032191,\qquad M=524288,
\]
the nonnegative representative bound is
\[
L(M-1)^2=283725419316756479
 <998244353\cdot1004535809
 =1002772198720536577.
\]
The common transform length $2^{21}$ divides both finite-field multiplicative periods. This bound, primitive-root checks and integer-type limits are part of the computation contract. Modular agreement alone would not justify CRT reconstruction without them.

\section{How to reproduce the record}
Paths in this section are relative to the extracted project root. The executable entry point is
\begin{lstlisting}[basicstyle=\ttfamily\small,breaklines=true]
python3 code/run_replays.py --output-root evidence/computation/reproduced
\end{lstlisting}
The output directory must be fresh. Individual suite names may precede the output option. The tested runtime and dependency versions are recorded in \path{evidence/BUILD_ENVIRONMENT.json}. Each execution records its actual arguments, inputs, result hashes and resource outcomes.

The full replay uses a $600$-second wall-clock bound per suite, a $570$-second per-process CPU cap, inherited one-core affinity and a two-million-byte combined log limit. Numerical-library thread caps are cooperative. No address-space cap is claimed. Exact arithmetic programs can consume substantial time or memory outside the recorded ranges.

The two readers rebuild with
\begin{lstlisting}[basicstyle=\ttfamily\small,breaklines=true]
python3 build.py --figures
python3 quality_check.py
python3 validate_fonts.py
\end{lstlisting}
Rendering creates immutable snapshots for inspection. Automated checks locate unresolved references, clipped text and font problems. Every final page in this edition also has an individual visual review record tied to the final PDF hash. A fresh rebuild may change PDF metadata, so its layout approval must be tied to its own output rather than copied from this edition.

\section{What the evidence does not claim}
No scientific replay failed. One inventory collector initially read the wrong schema field, then was corrected. Its failure and correction are recorded rather than hidden as a successful mathematical execution.

The recovered project contains written mathematics and exact programs, with no recovered Lehmer proof-assistant source project or compiled certificate. The formal inventory is an inventory and definition audit. It is not a Lean verification of the theorems in this book.

Each exact certificate proves its stated finite equality, rank or sign claim. Each range check ends at its stated bound. The universal deductions use the written proofs and the external contracts in Appendix~\ref{app:sources}. The all-prime nonvanishing statement remains open in this work.
